%=============================================================================!
% 16.8  DIFFUSION IN LAYERS                                                   !
%=============================================================================!
\section{Diffusion in layers}
\label{sec:ob-layers}

Lithium in graphite can move between the carbon sheets much more readily
than it can cross them. A single diffusion coefficient hides this
difference: we need the mean squared displacement in each direction.
After measuring it in the layered model solid, we connect diffusion to
hops between lattice sites and test the assumption that successive hops
are independent.

\subsection*{The diffusion tensor}

The MSD of \cref{eq:ob-einstein} is the sum of three parts, $\ens{\Delta
x^2} + \ens{\Delta y^2} + \ens{\Delta z^2}$, and each grows as $2D_{xx}t$,
$2D_{yy}t$ and $2D_{zz}t$ with its own coefficient. More generally the
products of the components define the diffusion tensor, $\ens{\Delta
r_\alpha\Delta r_\beta} \to 2D_{\alpha\beta}t$, a symmetric $3\times3$
matrix whose eigenvectors are the directions in which the spreading is
fastest and slowest (\cref{sec:os-modes}). In a layered solid with its
sheets at right angles to $z$, symmetry makes $z$ one of them and leaves
the two in the plane equal, $D_{xx} = D_{yy}$, and the diffusion
coefficient that a single number would give, a third of the trace,
mixes two very different motions.

\Cref{sec:pr-shape} built a model of such a solid: triangular sheets, each
bound strongly within itself and weakly to its neighbours, with 108 guest
atoms of lithium's mass grown into the hollows between them. Held at
\SI{150}{\kelvin} by CSVR in its relaxed cell for \SI{500}{\pico\second},
its guests' MSD along each axis, measured from the centre of mass, gives
from five blocks of \SI{100}{\pico\second} and fits over 5 to
\SI{50}{\pico\second}
\begin{equation*}
   D_{xx} = (1.16 \pm 0.11)\times10^{-4}\,\si{\angstrom\squared\per
   \femto\second} , \qquad
   D_{yy} = (1.02 \pm 0.05)\times10^{-4}\,\si{\angstrom\squared\per
   \femto\second} .
\end{equation*}
The two in the plane agree within their errors. Along $z$ the MSD is
\SI{0.113}{\angstrom\squared} at \SI{1}{\pico\second},
\SI{0.126}{\angstrom\squared} at 10 and \SI{0.132}{\angstrom\squared} at
\SI{50}{\pico\second} (\cref{fig:ob-layers}a), and no guest moves more
than \SI{1.73}{\angstrom} along $z$ in the whole run: the guests rattle
between their two sheets, and none crosses to the next gap,
\SI{3.40}{\angstrom} away. A single crossing would add
\SI{11.6}{\angstrom\squared} to one guest's square, and so
$\SI{11.6}{\angstrom\squared}/108$ to the mean, and one in the
\SI{500}{\pico\second} would have meant $D_{zz} \approx
\SI{1.1e-7}{\angstrom\squared\per\femto\second}$. This is the scale of one detectable crossing, not a statistical upper
bound on $D_{zz}$: a finite run can miss a rare event. It is a thousand
times smaller than the in-plane value. A straight line fitted to the
levelling MSD would only measure its slow approach to the plateau;
\cref{ex:ip-rate} shows how zero events give a confidence bound.

\subsection*{Hops between sites}

On the model surface of \cref{sec:en-surface} a lithium atom moves by
hops between hollows a distance $a = \SI{2.46}{\angstrom}$ apart, and
Chapter~13 counted them. If each hop goes to one of the six neighbouring
hollows at random, independently of the last, the atom performs the
random walk of \cref{sec:ob-msd} with steps of length $a$, and after $n$
hops its mean squared displacement is $na^2$. At a rate of $k_{\mathrm h}$ hops
per unit time, $n = k_{\mathrm h}t$, and comparing $k_{\mathrm h}ta^2$ with $2dDt$ for
$d = 2$,
\begin{equation*}
   D = \frac{k_{\mathrm h}a^2}{4} .
\end{equation*}
Counted hops thus give a diffusion coefficient without fitting an MSD,
and a count has the error bar of \cref{sec:ob-reactions}, $\sqrt n$. If
the hops have various lengths $\ell$ but still forget their direction,
the same sum gives $n\ens{\ell^2}$ after $n$ hops, and $D = k_{\mathrm
h}\ens{\ell^2}/4$.

The relation assumes that hops forget their direction. At
\SI{1000}{\kelvin}, with Langevin friction of \SI{1}{\per\pico\second}, 1000
atoms over \SI{200}{\pico\second} make \num{179680} hops, a rate of $k_{\mathrm h}
= (0.8984 \pm 0.0021)$ per atom and picosecond, Chapter~13's rate again;
$k_{\mathrm h}a^2/4$ is then \SI{13.6e-4}{\angstrom\squared\per\femto\second},
but the MSD over 2 to \SI{10}{\pico\second} gives $D =
\SI{92.3e-4}{\angstrom\squared\per\femto\second}$, 6.79 times more. The
ratio $f = 4D/k_{\mathrm h}a^2$, the correlation factor, exceeds 1 for two
reasons. An atom with enough energy to cross a bridge often keeps it long
enough to cross the next before the friction takes it away, so that 17\%
of the counted hops, from one hollow where the atom is seen to the next,
are longer than $a$, and their mean square is $1.96a^2$. And successive
hops tend to keep their direction: $4D/k_{\mathrm h}\ens{\ell^2}$, which would
be 1 for hops that forget it, is 3.46. Raising the friction to 3, 10 and
\SI{30}{\per\pico\second} drains the energy faster: $f$ falls to 3.24,
1.63 and 1.13, as the product of $\ens{\ell^2}/a^2$, 1.66, 1.29 and 1.07,
and the memory of direction, 1.96, 1.27 and 1.05 (\cref{fig:ob-layers}b),
while $k_{\mathrm h}$ stays near 0.90 until the highest friction, which slows
the hops themselves to 0.708 per picosecond, as Chapter~13 found. How large $f$ is in a real crystal
depends on how fast the lattice takes away an arriving atom's energy,
which only a measurement of both $D$ and $k_{\mathrm h}$ settles.

\begin{figure}[htb]
\centering
\includegraphics{ch16_observables/layers}
\caption{(a) The guests of the layered model solid at \SI{150}{\kelvin}:
their mean squared displacement along $x$, $y$ and $z$, measured from the
centre of mass. (b) Lithium on the model surface at \SI{1000}{\kelvin}:
the correlation factor $f = 4D/k_{\mathrm h}a^2$, with $D$ from the MSD and
$k_{\mathrm h}$ from counted hops (circles), and $4D/k_{\mathrm h}\ens{\ell^2}$,
with $\ens{\ell^2}$ the mean square length of a counted hop (squares),
against the Langevin friction. (c) The
vibrational density of states of the lithium atoms at 600 and
\SI{1000}{\kelvin} (friction \SI{1}{\per\pico\second}), with the
frequency of small vibrations in a hollow (dotted; \cref{sec:ob-vdos}).}
\label{fig:ob-layers}
\end{figure}

\begin{keyidea}
In an anisotropic solid the MSD is taken along each direction, giving the
diffusion tensor; in a layered solid the motion along the layers and
across them differ by orders of magnitude. Diffusion by hops between
sites a distance $a$ apart gives $D = fk_{\mathrm h}a^2/4$ in two dimensions,
with $f = 1$ when successive hops are independent and one site long, and
$f > 1$ when they carry the atom past the next site or keep their
direction.
\end{keyidea}

\notebook{16}{split the MSD of the guests into its three directions}

\begin{selfcheck}
On a simple cubic lattice of spacing $a$ in three dimensions, with
independent hops at a total rate $k_{\mathrm h}$, what is $D$?
\end{selfcheck}
